Để chuẩn bị cho phần này của workshop, bạn sẽ cần phải: + Triển khai
CloudFormation stack + Sửa đổi bảng định tuyến VPC.

Các thành phần này hoạt động cùng nhau để mô phỏng DNS forwarding và
name resolution.

\vspace{0.5em}\noindent\textbf{Triển khai CloudFormation
stack}\par

Mẫu CloudFormation sẽ tạo các dịch vụ bổ sung để hỗ trợ mô phỏng môi
trường truyền thống: + Một Route 53 Private Hosted Zone lưu trữ các bản
ghi Bí danh (Alias records) cho điểm cuối PrivateLink S3 + Một Route 53
Inbound Resolver endpoint cho phép ``VPC Cloud'' giải quyết các yêu cầu
resolve DNS gửi đến Private Hosted Zone + Một Route 53 Outbound Resolver
endpoint cho phép ``VPC On-prem'' chuyển tiếp các yêu cầu DNS cho S3
sang ``VPC Cloud''

\begin{figure}
\centering
\pandocbounded{\includegraphics[keepaspectratio,alt={route 53 diagram}]{5-Workshop/5.4-S3-onprem/route53.png}}
\caption{route 53 diagram}
\end{figure}

\begin{enumerate}
\def\labelenumi{\arabic{enumi}.}
\tightlist
\item
  Click link sau để mở
  \href{https://us-east-1.console.aws.amazon.com/cloudformation/home?region=us-east-1\#/stacks/quickcreate?templateURL=https://s3.amazonaws.com/reinvent-endpoints-builders-session/R53CF.yaml&stackName=PLOnpremSetup}{AWS
  CloudFormation console}. Mẫu yêu cầu sẽ được tải sẵn vào menu. Chấp
  nhận tất cả mặc định và nhấp vào Tạo stack.
\end{enumerate}

\begin{figure}
\centering
\pandocbounded{\includegraphics[keepaspectratio,alt={Create stack}]{5-Workshop/5.4-S3-onprem/create-stack.png}}
\caption{Create stack}
\end{figure}

\begin{figure}
\centering
\pandocbounded{\includegraphics[keepaspectratio,alt={Button}]{5-Workshop/5.4-S3-onprem/create-stack-button.png}}
\caption{Button}
\end{figure}

Có thể mất vài phút để triển khai stack hoàn tất. Bạn có thể tiếp tục
với bước tiếp theo mà không cần đợi quá trình triển khai kết thúc.

\vspace{0.5em}\noindent\textbf{Cập nhật bảng định tuyến private
on-premise}\par

Workshop này sử dụng StrongSwan VPN chạy trên EC2 instance để mô phỏng
khả năng kết nối giữa trung tâm dữ liệu truyền thống và môi trường cloud
AWS. Hầu hết các thành phần bắt buộc đều được cung cấp trước khi bạn bắt
đầu. Để hoàn tất cấu hình VPN, bạn sẽ sửa đổi bảng định tuyến ``VPC
on-prem'' để hướng lưu lượng đến cloud đi qua StrongSwan VPN instance.

\begin{enumerate}
\def\labelenumi{\arabic{enumi}.}
\item
  Mở Amazon EC2 console
\item
  Chọn instance tên infra-vpngw-test. Từ Details tab, copy Instance ID
  và paste vào text editor của bạn để sử dụng ở những bước tiếp theo
\end{enumerate}

\begin{figure}
\centering
\pandocbounded{\includegraphics[keepaspectratio,alt={ec2 id}]{5-Workshop/5.4-S3-onprem/ec2-onprem-id.png}}
\caption{ec2 id}
\end{figure}

\begin{enumerate}
\def\labelenumi{\arabic{enumi}.}
\setcounter{enumi}{2}
\item
  Đi đến VPC menu bằng cách gõ ``VPC'' vào Search box
\item
  Click vào Route Tables, chọn RT Private On-prem route table, chọn
  Routes tab, và click Edit Routes.
\end{enumerate}

\begin{figure}
\centering
\pandocbounded{\includegraphics[keepaspectratio,alt={rt}]{5-Workshop/5.4-S3-onprem/rt.png}}
\caption{rt}
\end{figure}

\begin{enumerate}
\def\labelenumi{\arabic{enumi}.}
\setcounter{enumi}{4}
\tightlist
\item
  Click Add route.
\end{enumerate}

\begin{itemize}
\tightlist
\item
  Destination: CIDR block của Cloud VPC
\item
  Target: ID của infra-vpngw-test instance (bạn đã lưu lại ở bước trên)
\end{itemize}

\begin{figure}
\centering
\pandocbounded{\includegraphics[keepaspectratio,alt={add route}]{5-Workshop/5.4-S3-onprem/add-route.png}}
\caption{add route}
\end{figure}

\begin{enumerate}
\def\labelenumi{\arabic{enumi}.}
\setcounter{enumi}{5}
\tightlist
\item
  Click Save changes
\end{enumerate}
